\section{敏感性分析}
\label{sec:sens}

\subsection{返航安全余量对最大安全载荷与组批的影响}

返航电量下限（安全余量比例 $\eta_g$）是贯穿四问的关键参数。这里单独引用代码包中的 \texttt{p1-1/solve\_p1\_1.py} 敏感性实验；它与正文主链 \texttt{q1\_21/q1\_22/q1\_23} 使用同一基础数据，但组批实现独立，因此不将其数值直接替换主链基准结果。实验在 $\eta=10\%,15\%,20\%,30\%,40\%$ 五个离散场景下重复求解。

独立实验显示：C 型最大安全载荷随安全余量增加持续下降；A、B 型在低安全余量下仍由质量上限控制，至 $\eta=30\%$ 才明显转为能量受限。当 $\eta=40\%$ 时场景不可行。平衡组批架次数在 $\eta=10\%,15\%,20\%$ 时均为 $19$，在 $\eta=30\%$ 时升至 $21$，因此不能把该实验概括为“$10\%\!\sim\!25\%$ 内始终为 $18$ 架次”。图~\ref{fig:sens} 给出该独立实验的真实载荷和架次响应。

\begin{figure}[htbp]
  \centering
  \includegraphics[width=0.96\textwidth]{p1-1_sensitivity.png}
  \caption{独立 p1-1 实验中的返航安全余量敏感性}
  \label{fig:sens}
\end{figure}

因此，主链在 $\eta=20\%$ 下得到的 $N=18$ 代表组批应理解为基准场景结论；独立敏感性实验只用于说明安全余量的方向性影响和不可行阈值，不用于声称主链三种算法在所有扰动场景下仍保持同一数值。

\subsection{其他参数与算法稳定性}

问题二、三的 NSGA-II 结果应以固定随机种子和汇总文件为准；本文不额外宣称未经汇总文件支持的“波动均在 $3\%$ 以内”。代码包中的 \texttt{p3-3pro} 在固定上游路线下优化访问顺序与中继点，得到 $29$ 条运输路线、$19$ 个中继架次、联合完工 $43993.76\,$s、能耗 $95.664\,$kWh、名义通信中断 $0$；其时限违反以累计秒数 $145177.47\,$s 记录，与主链的违反条数不是同一统计口径。另一个 \texttt{p3-4pro} 测试基准、附加 $5$ dB 和 $10$ dB 衰落，最坏情景通信中断为 $15631.94\,$s。两者均是互补增强结果，不与主解直接排名。
